\documentclass[11pt]{article}
\usepackage[utf8]{inputenc}
\usepackage[margin=1in]{geometry}
\usepackage{amsmath}
\usepackage{graphicx}
\usepackage{booktabs}
\usepackage{hyperref}
\usepackage[round]{natbib}
\usepackage{float}

\title{Potassium Regulates Axon--Oligodendrocyte Signaling and Metabolic Coupling in White Matter}
\author{Author A$^{1}$, Author B$^{1,2}$, Author C$^{2}$\\
\small $^{1}$Institute of Neuronal Cell Biology, Technical University of Munich, Germany\\
\small $^{2}$German Center for Neurodegenerative Diseases (DZNE), Munich, Germany}
\date{}

\begin{document}
\maketitle

\begin{abstract}
Oligodendrocytes provide metabolic support to myelinated axons through monocarboxylate transporters, but how they detect axonal activity and regulate this support is poorly understood. Here, we demonstrate that potassium released during high-frequency axonal firing activates inwardly rectifying Kir4.1 channels on mature oligodendrocytes, triggering a biphasic calcium response that scales with stimulation frequency ($p < 0.0001$, 5 vs.\ 30~mM K$^+$). This response is barium-sensitive and depends on extracellular calcium influx via reverse-mode Na$^+$/Ca$^{2+}$ exchange, while glutamatergic and purinergic signaling do not contribute (NBQX $p = 0.73$; PPADS not significant). Oligodendrocyte-specific Kir4.1 conditional knockout mice (cKO) developed delayed compound action potential recovery and axonal pathology by 7--8 months. Quantitative myelin proteomics revealed reduced GLUT1 and MCT1 transporter expression in 2--3 month-old cKO mice before onset of axonopathy. Using genetically encoded FRET sensors in optic nerve axons, we found lower resting lactate levels, reduced activity-induced lactate surges, decreased axonal glucose uptake, and dramatically reduced activity-driven glucose consumption in cKO axons (2.7 $\pm$ 0.6\%/min in controls vs.\ 0.1 $\pm$ 0.2\%/min in cKO; $p = 0.0002$). These results establish K$^+$-Kir4.1 signaling as a critical pathway through which oligodendrocytes sense axonal activity and regulate metabolic support, with disruption leading to progressive axonal pathology.
\end{abstract}

\section{Introduction}

Myelinating oligodendrocytes (OLs) provide essential metabolic support to axons by supplying lactate and pyruvate through monocarboxylate transporters (MCTs) in the myelin sheath \citep{funfschilling2012glycolytic, lee2012oligodendroglia, nave2010myelination}. Loss of this metabolic coupling leads to axonal degeneration even when myelin structure remains intact \citep{philips2017oligodendroglia, simons2014oligodendrocytes}. However, the molecular mechanisms by which OLs detect axonal activity and adjust metabolic support remain largely unknown.

The inwardly rectifying potassium channel Kir4.1 is highly expressed in mature OLs, where it is thought to participate in extracellular K$^+$ clearance during neuronal activity \citep{neusch2001kir41, larson2018oligodendrocyte}. During action potential propagation, K$^+$ accumulates in the periaxonal space beneath the myelin sheath, creating a local signal of axonal activity. We hypothesised that this K$^+$ elevation activates Kir4.1 channels on OLs, triggering intracellular signaling cascades that regulate metabolic transporter expression and metabolite supply.

Here, we combine ex vivo optic nerve electrophysiology, two-photon calcium imaging in genetically labelled OLs, OL-specific Kir4.1 conditional knockout, quantitative myelin proteomics, and FRET-based metabolite imaging in axons to establish the K$^+$--Kir4.1 pathway as a critical mediator of activity-dependent axon--OL metabolic coupling.

\section{Results}

\subsection{Frequency-Dependent Calcium Responses in Oligodendrocytes}

Ex vivo optic nerves from PLP-CreERT;RCL-GCaMP6s mice were stimulated at 10, 25, and 50~Hz while recording OL calcium dynamics (Table~\ref{tab:calcium}).

\begin{table}[H]
\centering
\caption{Stimulus-evoked OL calcium responses by frequency ($n = 108$ cells, $N = 8$ nerves).}
\label{tab:calcium}
\begin{tabular}{@{}lccc@{}}
\toprule
Frequency (Hz) & Ca$^{2+}$ AUC (a.u.) & CAP Amplitude Change (\%) & Response Pattern \\
\midrule
10 & $1.42 \pm 0.31$ & $-8.3$ & Biphasic \\
25 & $2.87 \pm 0.48$ & $-16.7$ & Biphasic \\
50 & $4.61 \pm 0.72$ & $-24.1$ & Biphasic \\
\bottomrule
\end{tabular}
\end{table}

Calcium responses were biphasic (initial rise followed by transient undershoot) and scaled with stimulation frequency. TTX (1~$\mu$M) abolished the response entirely, confirming dependence on action potential propagation. Removal of extracellular calcium strongly diminished the response, indicating a calcium influx mechanism.

\subsection{K$^+$-Evoked Calcium Responses}

Bath application of elevated K$^+$ mimicked stimulus-evoked responses in a dose-dependent manner (Table~\ref{tab:potassium}).

\begin{table}[H]
\centering
\caption{K$^+$-evoked OL calcium responses (one-way ANOVA with Tukey's test, $n = 35$--$57$ cells, $N = 4$--$5$ nerves).}
\label{tab:potassium}
\begin{tabular}{@{}lccl@{}}
\toprule
$\Delta$[K$^+$] (mM) & Ca$^{2+}$ AUC (a.u.) & Comparison & $p$-value \\
\midrule
5 & $0.87 \pm 0.22$ & 5 vs.\ 10 & $\mathbf{0.0048}$ \\
10 & $1.94 \pm 0.38$ & 5 vs.\ 30 & $\mathbf{< 0.0001}$ \\
30 & $4.23 \pm 0.67$ & 10 vs.\ 30 & $\mathbf{< 0.0001}$ \\
\bottomrule
\end{tabular}
\end{table}

\subsection{Pharmacological Dissection}

Pharmacological experiments identified Kir channels and reverse-mode NCX as mediators, while glutamatergic and purinergic receptors were not involved (Table~\ref{tab:pharma}).

\begin{table}[H]
\centering
\caption{Pharmacological dissection of OL calcium response mechanisms.}
\label{tab:pharma}
\begin{tabular}{@{}llccl@{}}
\toprule
Agent & Target & Effect on OL Ca$^{2+}$ & $p$-value & Test \\
\midrule
Barium & Kir channels & Diminished & $< 0.01$ & Paired $t$-test \\
NBQX (50~$\mu$M) & AMPA/KA & No effect & $0.7277$ & Paired $t$-test \\
PPADS (100~$\mu$M) & P2 receptors & No effect & $0.42$ & Paired $t$-test \\
KB-R7943 & NCX (reverse) & Partial reduction & $< 0.05$ & Paired $t$-test \\
0~mM Ca$^{2+}$ & Ca$^{2+}$ influx & Strong reduction & $< 0.001$ & Paired $t$-test \\
\bottomrule
\end{tabular}
\end{table}

Barium did not affect axonal calcium surges, confirming OL-specific Kir channel involvement.

\subsection{Kir4.1 cKO: CAP Recovery and K$^+$ Clearance}

OL-specific Kir4.1 cKO mice showed normal baseline CAP properties but impaired recovery after high-frequency stimulation (Table~\ref{tab:cap}).

\begin{table}[H]
\centering
\caption{CAP properties in ctrl ($n = 8$) vs.\ cKO ($n = 9$) optic nerves. ANOVA with Holm-Sidak correction.}
\label{tab:cap}
\begin{tabular}{@{}lccc@{}}
\toprule
Measure & ctrl & cKO & $p$-value \\
\midrule
Peak 1 latency (ms) & $1.23 \pm 0.08$ & $1.24 \pm 0.09$ & $0.764$ \\
Peak 2 latency (ms) & $2.41 \pm 0.12$ & $2.41 \pm 0.14$ & $0.996$ \\
Peak 3 latency (ms) & $3.87 \pm 0.21$ & $3.89 \pm 0.24$ & $0.927$ \\
Post-stim recovery & Normal & Delayed & $\mathbf{< 0.01}$ \\
\bottomrule
\end{tabular}
\end{table}

\subsection{Myelin Proteomics}

TMT-labeled LC--MS/MS on optic nerves from 2.5-month-old mice revealed reduced metabolic transporter expression in cKO myelin (Table~\ref{tab:proteomics}).

\begin{table}[H]
\centering
\caption{Key protein changes in cKO vs.\ ctrl myelin ($n = 4$ per group).}
\label{tab:proteomics}
\begin{tabular}{@{}lccc@{}}
\toprule
Protein & Gene & log$_2$ Fold Change & Significant? \\
\midrule
GLUT1 & \textit{Slc2a1} & $-0.47$ & \textbf{Yes} \\
MCT1 & \textit{Slc16a1} & $-0.38$ & \textbf{Yes} \\
Myelin basic protein & \textit{Mbp} & $+0.04$ & No \\
Proteolipid protein & \textit{Plp1} & $-0.08$ & No \\
\bottomrule
\end{tabular}
\end{table}

GLUT1 and MCT1 were significantly downregulated before any detectable axonopathy, suggesting that metabolic disruption precedes structural damage.

\subsection{Axonal Metabolite Dynamics}

FRET sensors delivered via intravitreal AAV injection revealed metabolic deficits in cKO axons (Table~\ref{tab:metabolites}).

\begin{table}[H]
\centering
\caption{Axonal metabolite measurements in ctrl vs.\ cKO optic nerves.}
\label{tab:metabolites}
\begin{tabular}{@{}llccl@{}}
\toprule
Sensor & Measure & ctrl & cKO & $p$-value \\
\midrule
Laconic & Resting lactate (ratio) & $1.14 \pm 0.08$ & $0.97 \pm 0.06$ & $\mathbf{< 0.01}$ \\
Laconic & Activity-induced surge & Present & Reduced & $\mathbf{< 0.05}$ \\
FLIIP & Resting glucose (ratio) & $0.82 \pm 0.05$ & $0.71 \pm 0.04$ & $\mathbf{< 0.01}$ \\
FLIIP & Glucose consumption (\%/min) & $2.7 \pm 0.6$ & $0.1 \pm 0.2$ & $\mathbf{0.0002}$ \\
ATeam & ATP (ratio change) & $-0.12 \pm 0.03$ & $-0.14 \pm 0.04$ & $0.18$ \\
\bottomrule
\end{tabular}
\end{table}

The most striking finding was the near-abolition of activity-driven glucose consumption in cKO axons (0.1 vs.\ 2.7\%/min), revealing that OL Kir4.1 is essential for regulating axonal glucose metabolism during neuronal activity.

\section{Discussion}

This study establishes K$^+$--Kir4.1 signaling as a critical pathway for activity-dependent axon--OL metabolic coupling. The proposed mechanism is that K$^+$ released during axonal spiking activates Kir4.1 channels on OLs, causing membrane depolarisation and Ca$^{2+}$ influx partly via reverse-mode NCX. This calcium signal adjusts OL metabolic support by regulating GLUT1 and MCT1 expression, lactate supply, and glucose transport. Loss of Kir4.1 disrupts this signaling, producing early metabolic deficits that precede and likely cause progressive axonal degeneration.

The frequency-dependent scaling of the OL calcium response provides a plausible mechanism for activity-dependent metabolic matching: more active axons produce more K$^+$, eliciting larger OL responses and greater metabolic support \citep{nave2010myelination, simons2014oligodendrocytes}. The insensitivity to NBQX and PPADS rules out glutamatergic and purinergic transmission as primary mediators in mature white matter, distinguishing this pathway from signaling to oligodendrocyte precursor cells \citep{bergles2000glutamatergic}.

The clinical relevance of these findings extends to multiple neurological conditions. Kir4.1 autoantibodies have been identified in multiple sclerosis, and Kir4.1 dysfunction may contribute to axonal pathology in various demyelinating and neurodegenerative conditions \citep{larson2018oligodendrocyte}. Understanding the molecular basis of axon--OL metabolic coupling may reveal new therapeutic targets for protecting axonal integrity.

\section{Methods}

\subsection{Animals}
PLP-CreERT;RCL-GCaMP6s mice for OL calcium imaging. Kir4.1$^{fl/fl}$;MOGiCre for OL-specific Kir4.1 knockout. AAV-ATeam1.03, AAV-Laconic, and AAV-FLIIP for axonal metabolite imaging via intravitreal injection.

\subsection{Ex Vivo Optic Nerve Preparation}
Optic nerves were dissected and maintained in oxygenated ACSF at 37$^\circ$C. Compound action potentials were recorded using suction electrodes. Two-photon imaging captured GCaMP6s and FRET sensor signals. Stimulation used 30-second trains at 10, 25, or 50~Hz.

\subsection{Pharmacology}
TTX (1~$\mu$M), barium chloride, NBQX (50~$\mu$M), PPADS (100~$\mu$M), KB-R7943, and 0~mM extracellular Ca$^{2+}$ were applied to identify signaling mechanisms.

\subsection{Proteomics}
TMT-labeled LC--MS/MS on optic nerves from 2.5-month-old cKO ($n = 4$) and ctrl ($n = 4$) mice. Proteins ranked by log$_2$ fold change.

\subsection{Statistics}
One-way ANOVA with Tukey's or Holm-Sidak corrections for multiple comparisons. Paired and unpaired $t$-tests for two-group comparisons.

\section{Conclusion}

Potassium-mediated activation of oligodendroglial Kir4.1 channels constitutes a previously unknown signaling pathway that couples axonal activity to metabolic support in white matter. Loss of this pathway produces early downregulation of metabolic transporters, impaired lactate and glucose supply to axons, and eventual axonal degeneration, establishing K$^+$--Kir4.1 as an essential mediator of axon--oligodendrocyte metabolic homeostasis.

\bibliographystyle{plainnat}
\bibliography{references}

\end{document}
